% GENERATED FILE. Edit canonical lessons, then run npm run generate:books.
% canonical-source-hash: fnv1a64:45c0235e617a0ee3
% canonical-lessons: PA-C193-ordinal-return, PA-C193-pahile, PA-C193-pahile-review, PA-C193-duje, PA-C193-duje-review, PA-C193-ordinal-contrast, PA-C193-ordinal-frames, PA-C193-payoff

\chapter{First and Second in a Room}
\label{ch:pa-a2-ordinal-e}
\hlchaptermodality{193}

\begin{chapteropening}
\textbf{By the end of this chapter:} \emph{I can contrast first and second across a direct room phrase and the same masculine oblique location frame.}

Contrast first and second across direct and oblique room frames through listening, speaking, reading, and writing.
\end{chapteropening}

\section[\texorpdfstring{\pa{ਪਹਿਲਾ} · \pa{ਦੂਜਾ} … (pahilaa; duujaa …)}{ਪਹਿਲਾ · ਦੂਜਾ … (pahilaa; duujaa …)}]{Return to first, second, and room}

\label{lesson:PA-C193-ordinal-return}



\begin{quote}
\emph{Hear:} \emph{pahilaa; duujaa}; YOU SAY: \textbf{\pa{ਪਹਿਲਾ}; \pa{ਦੂਜਾ}}
\end{quote}

\begin{grammarlens}[title={Grammar Lens}]
Keep two familiar room frames separate: \textbf{\pa{ਕਮਰਾ}} on its own and \textbf{\pa{ਕਮਰੇ} \pa{ਵਿੱਚ}} for “in the room.” No new ordinal form enters this opener.
\end{grammarlens}

\subsection*{Your turn}
\begin{itemize}
\raggedright
  \item \emph{Hear:} \emph{kamraa}; YOU SAY: \textbf{\pa{ਕਮਰਾ}}
  \item \emph{Hear:} \emph{kamre vich}; YOU SAY: \textbf{\pa{ਕਮਰੇ} \pa{ਵਿੱਚ}}
  \item \emph{Hear:} \emph{pahilaa; duujaa}; YOU SAY: \textbf{\pa{ਪਹਿਲਾ}; \pa{ਦੂਜਾ}}
\end{itemize}

\begin{tcolorbox}[breakable,colback=teal!4,colframe=teal!35!black,title={Before you move on}]
Say “first,” “second,” “room,” and “in the room.” This opener owns no new form.
\end{tcolorbox}
\section[\texorpdfstring{\pa{ਪਹਿਲੇ} (pahile)}{ਪਹਿਲੇ (pahile)}]{Change first inside one familiar room frame}

\label{lesson:PA-C193-pahile}



\begin{quote}
Retrieve \textbf{\pa{ਪਹਿਲਾ}}, then \textbf{\pa{ਕਮਰਾ}}, then \textbf{\pa{ਕਮਰੇ} \pa{ਵਿੱਚ}}.
\end{quote}

\begin{grammarlens}[title={Grammar Lens: direct owner}]
\begin{quote}
\textbf{\pa{ਪਹਿਲਾ} \pa{ਕਮਰਾ}} — \emph{pahilā kamrā} — the first room
\end{quote}

\begin{quote}
\textbf{\pa{ਪਹਿਲੇ} \pa{ਕਮਰੇ} \pa{ਵਿੱਚ}} — \emph{pahile kamre vicc} — in the first room
\end{quote}

The familiar direct form has the final \emph{-aa} sound. In this masculine singular phrase, the noun is oblique before \textbf{\pa{ਵਿੱਚ}}, and the ordinal takes final \emph{-e}.
\end{grammarlens}

\begin{grammarlens}[title={Grammar Lens: source boundary}]
Gill and Gleason's masculine adjective table contrasts the singular direct and singular oblique patterns. Here that contrast is \textbf{\pa{ਪਹਿਲਾ}} beside \textbf{\pa{ਪਹਿਲੇ}}. Their following numeral section lists first \textbf{\pa{ਪਹਿਲਾ}} and second \textbf{\pa{ਦੂਜਾ}} among Punjabi ordinals. This chapter follows these two ordinals, not a claim that every Punjabi adjective changes this way. Source: Harjeet Singh Gill and Henry A. Gleason, Jr., \emph{A Reference Grammar of Punjabi} (1969), pp. 71–73, \href{https://apnaorg.com/books/english/a-reference-grammar-of-punjabi/a-reference-grammar-of-punjabi.pdf}{adjectives and numerals}.
\end{grammarlens}

\subsection*{Your turn}
\begin{itemize}
\raggedright
  \item \emph{Hear:} \emph{pahilaa kamraa}; YOU SAY: \textbf{\pa{ਪਹਿਲਾ} \pa{ਕਮਰਾ}}
  \item \emph{Hear:} \emph{pahile kamre vich}; YOU SAY: \textbf{\pa{ਪਹਿਲੇ} \pa{ਕਮਰੇ} \pa{ਵਿੱਚ}}
  \item \emph{Read it:} \textbf{\pa{ਪਹਿਲਾ} \pa{ਕਮਰਾ} · \pa{ਪਹਿਲੇ} \pa{ਕਮਰੇ} \pa{ਵਿੱਚ}}
  \item \emph{Write it:} copy \textbf{\pa{ਪਹਿਲੇ}} once, then compare it with \textbf{\pa{ਪਹਿਲਾ}}
\end{itemize}

\begin{tcolorbox}[breakable,colback=teal!4,colframe=teal!35!black,title={Before you move on}]
Produce “in the first room.”
\end{tcolorbox}
\section[\texorpdfstring{\pa{ਪਹਿਲਾ} · \pa{ਪਹਿਲੇ} (pahilaa; pahile)}{ਪਹਿਲਾ · ਪਹਿਲੇ (pahilaa; pahile)}]{Retrieve first in two small frames}

\label{lesson:PA-C193-pahile-review}



\begin{quote}
Say the familiar direct form \textbf{\pa{ਪਹਿਲਾ}}.
\end{quote}

\begin{grammarlens}[title={Grammar Lens}]
Keep one contrast: \textbf{\pa{ਪਹਿਲਾ} \pa{ਕਮਰਾ}}; \textbf{\pa{ਪਹਿਲੇ} \pa{ਕਮਰੇ} \pa{ਵਿੱਚ}}.
\end{grammarlens}

\subsection*{Your turn}
\begin{itemize}
\raggedright
  \item \emph{Hear:} \emph{pahilaa kamraa; pahile kamre vich}; YOU SAY: \textbf{\pa{ਪਹਿਲਾ} \pa{ਕਮਰਾ}; \pa{ਪਹਿਲੇ} \pa{ਕਮਰੇ} \pa{ਵਿੱਚ}}
  \item \emph{Read it:} \textbf{\pa{ਪਹਿਲੇ} \pa{ਕਮਰੇ} \pa{ਵਿੱਚ}}
  \item \emph{Write it:} cover and write \textbf{\pa{ਪਹਿਲੇ}}
\end{itemize}

\begin{tcolorbox}[breakable,colback=teal!4,colframe=teal!35!black,title={Before you move on}]
Produce the form of “first” used in the \textbf{\pa{ਕਮਰੇ} \pa{ਵਿੱਚ}} frame.
\end{tcolorbox}
\section[\texorpdfstring{\pa{ਦੂਜੇ} (duuje)}{ਦੂਜੇ (duuje)}]{Apply the same bounded change to second}

\label{lesson:PA-C193-duje}



\begin{quote}
Retrieve \textbf{\pa{ਦੂਜਾ}}, then the room-frame change you just practised.
\end{quote}

\begin{grammarlens}[title={Grammar Lens: direct owner}]
\begin{quote}
\textbf{\pa{ਦੂਜਾ} \pa{ਕਮਰਾ}} — \emph{dūjā kamrā} — the second room
\end{quote}

\begin{quote}
\textbf{\pa{ਦੂਜੇ} \pa{ਕਮਰੇ} \pa{ਵਿੱਚ}} — \emph{dūje kamre vicc} — in the second room
\end{quote}

The boundary stays narrow: the familiar direct ordinal changes from final \emph{-aa} to \emph{-e} in this same masculine singular oblique room frame.
\end{grammarlens}

\subsection*{Your turn}
\begin{itemize}
\raggedright
  \item \emph{Hear:} \emph{duujaa kamraa}; YOU SAY: \textbf{\pa{ਦੂਜਾ} \pa{ਕਮਰਾ}}
  \item \emph{Hear:} \emph{duuje kamre vich}; YOU SAY: \textbf{\pa{ਦੂਜੇ} \pa{ਕਮਰੇ} \pa{ਵਿੱਚ}}
  \item \emph{Read it:} \textbf{\pa{ਦੂਜਾ} \pa{ਕਮਰਾ} · \pa{ਦੂਜੇ} \pa{ਕਮਰੇ} \pa{ਵਿੱਚ}}
  \item \emph{Write it:} copy \textbf{\pa{ਦੂਜੇ}} once, then compare it with \textbf{\pa{ਦੂਜਾ}}
\end{itemize}

\begin{tcolorbox}[breakable,colback=teal!4,colframe=teal!35!black,title={Before you move on}]
Produce “in the second room.”
\end{tcolorbox}
\section[\texorpdfstring{\pa{ਦੂਜਾ} · \pa{ਦੂਜੇ} (duujaa; duuje)}{ਦੂਜਾ · ਦੂਜੇ (duujaa; duuje)}]{Retrieve second in two small frames}

\label{lesson:PA-C193-duje-review}



\begin{quote}
Say the familiar direct form \textbf{\pa{ਦੂਜਾ}}.
\end{quote}

\begin{grammarlens}[title={Grammar Lens}]
Keep one contrast: \textbf{\pa{ਦੂਜਾ} \pa{ਕਮਰਾ}}; \textbf{\pa{ਦੂਜੇ} \pa{ਕਮਰੇ} \pa{ਵਿੱਚ}}.
\end{grammarlens}

\subsection*{Your turn}
\begin{itemize}
\raggedright
  \item \emph{Hear:} \emph{duujaa kamraa; duuje kamre vich}; YOU SAY: \textbf{\pa{ਦੂਜਾ} \pa{ਕਮਰਾ}; \pa{ਦੂਜੇ} \pa{ਕਮਰੇ} \pa{ਵਿੱਚ}}
  \item \emph{Read it:} \textbf{\pa{ਦੂਜੇ} \pa{ਕਮਰੇ} \pa{ਵਿੱਚ}}
  \item \emph{Write it:} cover and write \textbf{\pa{ਦੂਜੇ}}
\end{itemize}

\begin{tcolorbox}[breakable,colback=teal!4,colframe=teal!35!black,title={Before you move on}]
Produce the form of “second” used in the \textbf{\pa{ਕਮਰੇ} \pa{ਵਿੱਚ}} frame.
\end{tcolorbox}
\section[\texorpdfstring{\pa{ਪਹਿਲਾ} · \pa{ਪਹਿਲੇ} … (pahilaa; pahile …)}{ਪਹਿਲਾ · ਪਹਿਲੇ … (pahilaa; pahile …)}]{Put both ordinal contrasts side by side}

\label{lesson:PA-C193-ordinal-contrast}



\begin{quote}
Retrieve \textbf{\pa{ਪਹਿਲਾ}} and \textbf{\pa{ਦੂਜਾ}} without a following noun.
\end{quote}

\begin{grammarlens}[title={Grammar Lens}]
\noindent
\begin{tabularx}{\linewidth}{@{}>{\raggedright\arraybackslash}X>{\raggedright\arraybackslash}X@{}}
\toprule
\textbf{Direct room frame} & \textbf{Oblique room frame} \\
\midrule
\textbf{\pa{ਪਹਿਲਾ} \pa{ਕਮਰਾ}} & \textbf{\pa{ਪਹਿਲੇ} \pa{ਕਮਰੇ} \pa{ਵਿੱਚ}} \\
\textbf{\pa{ਦੂਜਾ} \pa{ਕਮਰਾ}} & \textbf{\pa{ਦੂਜੇ} \pa{ਕਮਰੇ} \pa{ਵਿੱਚ}} \\
\bottomrule
\end{tabularx}

Read down one column before crossing to the other. The ending change belongs to this bounded masculine room contrast.
\end{grammarlens}

\subsection*{Your turn}
\begin{itemize}
\raggedright
  \item \emph{Hear:} \emph{pahilaa kamraa; pahile kamre vich}; YOU SAY: \textbf{\pa{ਪਹਿਲਾ} \pa{ਕਮਰਾ}; \pa{ਪਹਿਲੇ} \pa{ਕਮਰੇ} \pa{ਵਿੱਚ}}
  \item \emph{Hear:} \emph{duujaa kamraa; duuje kamre vich}; YOU SAY: \textbf{\pa{ਦੂਜਾ} \pa{ਕਮਰਾ}; \pa{ਦੂਜੇ} \pa{ਕਮਰੇ} \pa{ਵਿੱਚ}}
  \item \emph{Write it:} cover the table and write the two location-frame ordinal forms
\end{itemize}

\begin{tcolorbox}[breakable,colback=teal!4,colframe=teal!35!black,title={Before you move on}]
Say both “in the room” ordinal phrases.
\end{tcolorbox}
\section[\texorpdfstring{\pa{ਪਹਿਲੇ} \pa{ਕਮਰੇ} \pa{ਵਿੱਚ} … (pahile kamre …)}{ਪਹਿਲੇ ਕਮਰੇ ਵਿੱਚ … (pahile kamre …)}]{Choose first or second in the same location frame}

\label{lesson:PA-C193-ordinal-frames}



\begin{quote}
Produce the two new forms without adding the room phrase.
\end{quote}

\subsection*{Your turn}
\begin{itemize}
\raggedright
  \item \emph{Hear:} \emph{pahile kamre vich}; YOU SAY: \textbf{\pa{ਪਹਿਲੇ} \pa{ਕਮਰੇ} \pa{ਵਿੱਚ}}
  \item \emph{Hear:} \emph{duuje kamre vich}; YOU SAY: \textbf{\pa{ਦੂਜੇ} \pa{ਕਮਰੇ} \pa{ਵਿੱਚ}}
  \item \emph{Read it:} \textbf{\pa{ਪਹਿਲੇ} \pa{ਕਮਰੇ} \pa{ਵਿੱਚ} · \pa{ਦੂਜੇ} \pa{ਕਮਰੇ} \pa{ਵਿੱਚ}}
  \item \emph{Write it:} cover the line and write one phrase, then the other
\end{itemize}

\begin{grammarlens}[title={Grammar Lens}]
The choice changes the ordinal meaning. The familiar \textbf{\pa{ਕਮਰੇ} \pa{ਵਿੱਚ}} frame and the final \emph{-e} stay stable.
\end{grammarlens}

\begin{tcolorbox}[breakable,colback=teal!4,colframe=teal!35!black,title={Before you move on}]
Say “in the first room,” then “in the second room.”
\end{tcolorbox}
\section[\texorpdfstring{\pa{ਪਹਿਲਾ} \pa{ਕਮਰਾ} … (pahilaa kamraa …)}{ਪਹਿਲਾ ਕਮਰਾ … (pahilaa kamraa …)}]{A bounded first-and-second payoff}

\label{lesson:PA-C193-payoff}



\begin{quote}
Recall both direct ordinals, then both forms used in the familiar location frame.
\end{quote}

\subsection*{Your turn: four skills}
\begin{itemize}
\raggedright
  \item \emph{Hear:} \emph{pahilaa kamraa; pahile kamre vich; duujaa kamraa; duuje kamre vich}; YOU SAY: \textbf{\pa{ਪਹਿਲਾ} \pa{ਕਮਰਾ}; \pa{ਪਹਿਲੇ} \pa{ਕਮਰੇ} \pa{ਵਿੱਚ}; \pa{ਦੂਜਾ} \pa{ਕਮਰਾ}; \pa{ਦੂਜੇ} \pa{ਕਮਰੇ} \pa{ਵਿੱਚ}}
  \item \emph{Read it:} \textbf{\pa{ਪਹਿਲਾ} \pa{ਕਮਰਾ} · \pa{ਪਹਿਲੇ} \pa{ਕਮਰੇ} \pa{ਵਿੱਚ} · \pa{ਦੂਜਾ} \pa{ਕਮਰਾ} · \pa{ਦੂਜੇ} \pa{ਕਮਰੇ} \pa{ਵਿੱਚ}}
  \item \emph{Write it:} cover the line and write the two location phrases
\end{itemize}

\begin{tcolorbox}[breakable,colback=teal!4,colframe=teal!35!black,title={Before you move on}]
Produce both direct room phrases, then both location phrases.

This is the sixteenth prerequisite tranche exposed by Punjabi A2 mock drafting. It does not unblock Punjabi A2 mock 1 by itself. It is not an A2 mock, not A2 readiness or attainment, and not external or human validation.
\end{tcolorbox}
